\documentclass[12pt]{article}
\usepackage[UTF8]{inputenc}
\usepackage{xeCJK}
\setCJKmainfont{Sarabun}
\usepackage{enumitem}
\usepackage{geometry}
\geometry{a4paper, margin=2.5cm}
\usepackage{setspace}
\onehalfspacing

\begin{document}

\begin{center}
{\Large\textbf{ทบทวน (Review) (เพิ่มเติม)}}\\[6pt]
{\normalsize แบบฝึกหัดเพิ่มเติม}\\[6pt]
{\normalsize วิชาฟิสิกส์ ม.ปลาย บทที่ 30}\\[4pt]
{\footnotesize ทั้งหมด 30 ข้อ}
\end{center}
\hrule
\bigskip

\section*{แบบฝึกหัดเพิ่มเติม}
\begin{enumerate}[label=\arabic*., itemsep=0.5\baselineskip, topsep=0.5\baselineskip, leftmargin=*]
\item กฎข้อที่ 2 ของนิวตันระบุว่า F = ma หมายความว่าอะไร?
\begin{enumerate}[label=\alph*)., itemsep=0.3\baselineskip, topsep=0.3\baselineskip]
\item ก\ 
\begin{minipage}[t]{0.9\linewidth}
\raggedright แรง = มวล × ความเร่ง
\end{minipage}
\item ข\ 
\begin{minipage}[t]{0.9\linewidth}
\raggedright แรง = มวล × ความเร็ว
\end{minipage}
\item ค\ 
\begin{minipage}[t]{0.9\linewidth}
\raggedright แรง = มวล × ระยะทาง
\end{minipage}
\item ง\ 
\begin{minipage}[t]{0.9\linewidth}
\raggedright แรง = มวล/ความเร่ง
\end{minipage}
\end{enumerate}
\vspace{-0.3\baselineskip}
\begin{small}
\textit{คำตอบ: ก}
\end{small}

\item วัตถุอยู่ในสมดุลสถิติ หมายความว่าอะไร?
\begin{enumerate}[label=\alph*)., itemsep=0.3\baselineskip, topsep=0.3\baselineskip]
\item ก\ 
\begin{minipage}[t]{0.9\linewidth}
\raggedright ไม่มีแรงกระทำ
\end{minipage}
\item ข\ 
\begin{minipage}[t]{0.9\linewidth}
\raggedright ผลรวมแรง = 0 และผลรวมโมเมนต์ = 0
\end{minipage}
\item ค\ 
\begin{minipage}[t]{0.9\linewidth}
\raggedright ไม่เคลื่อนที่และไม่หมุน
\end{minipage}
\item ง\ 
\begin{minipage}[t]{0.9\linewidth}
\raggedright ข้อ ข และ ค ถูก
\end{minipage}
\end{enumerate}
\vspace{-0.3\baselineskip}
\begin{small}
\textit{คำตอบ: ง}
\end{small}

\item กฎของคูลอมบ์: F = kq₁q₂/r² ค่า k มีค่าเท่าใดในหน่วย SI?
\begin{enumerate}[label=\alph*)., itemsep=0.3\baselineskip, topsep=0.3\baselineskip]
\item ก\ 
\begin{minipage}[t]{0.9\linewidth}
\raggedright 9 × 10⁹ N·m²/C²
\end{minipage}
\item ข\ 
\begin{minipage}[t]{0.9\linewidth}
\raggedright 9 × 10⁶ N·m²/C²
\end{minipage}
\item ค\ 
\begin{minipage}[t]{0.9\linewidth}
\raggedright 8.85 × 10⁻¹² N·m²/C²
\end{minipage}
\item ง\ 
\begin{minipage}[t]{0.9\linewidth}
\raggedright 6.67 × 10⁻¹¹ N·m²/C²
\end{minipage}
\end{enumerate}
\vspace{-0.3\baselineskip}
\begin{small}
\textit{คำตอบ: ก}
\end{small}

\item คลื่นตามขวาง (Transverse wave) คือคลื่นที่มีทิศสั่นสะเทือน...?
\begin{enumerate}[label=\alph*)., itemsep=0.3\baselineskip, topsep=0.3\baselineskip]
\item ก\ 
\begin{minipage}[t]{0.9\linewidth}
\raggedright ขนานกับทิศการเคลื่อนที่ของคลื่น
\end{minipage}
\item ข\ 
\begin{minipage}[t]{0.9\linewidth}
\raggedright ตั้งฉากกับทิศการเคลื่อนที่ของคลื่น
\end{minipage}
\item ค\ 
\begin{minipage}[t]{0.9\linewidth}
\raggedright เป็นวงกลม
\end{minipage}
\item ง\ 
\begin{minipage}[t]{0.9\linewidth}
\raggedright ไม่มีทิศทาง
\end{minipage}
\end{enumerate}
\vspace{-0.3\baselineskip}
\begin{small}
\textit{คำตอบ: ข}
\end{small}

\item แบบจำลองอะตอมของ Rutherford มีสมมติฐานว่าอะตอมประกอบด้วย...?
\begin{enumerate}[label=\alph*)., itemsep=0.3\baselineskip, topsep=0.3\baselineskip]
\item ก\ 
\begin{minipage}[t]{0.9\linewidth}
\raggedright นิวเคลียสเล็กๆ ที่มีโปรตอน ล้อมรอบด้วยอิเล็กตรอน
\end{minipage}
\item ข\ 
\begin{minipage}[t]{0.9\linewidth}
\raggedright อะตอมเป็นทรงกลมทึบ
\end{minipage}
\item ค\ 
\begin{minipage}[t]{0.9\linewidth}
\raggedright อิเล็กตรอนอยู่ในวงโคจรแบบคงที่
\end{minipage}
\item ง\ 
\begin{minipage}[t]{0.9\linewidth}
\raggedright ไม่มีนิวเคลียส
\end{minipage}
\end{enumerate}
\vspace{-0.3\baselineskip}
\begin{small}
\textit{คำตอบ: ก}
\end{small}

\item คาบ (Period) ของลูกตุ้มอย่างง่าย T = 2π√(L/g) โดย L คือ...?
\begin{enumerate}[label=\alph*)., itemsep=0.3\baselineskip, topsep=0.3\baselineskip]
\item ก\ 
\begin{minipage}[t]{0.9\linewidth}
\raggedright มวลของลูกตุ้ม
\end{minipage}
\item ข\ 
\begin{minipage}[t]{0.9\linewidth}
\raggedright ความยาวเชือก
\end{minipage}
\item ค\ 
\begin{minipage}[t]{0.9\linewidth}
\raggedright ความถี่เชิงมุม
\end{minipage}
\item ง\ 
\begin{minipage}[t]{0.9\linewidth}
\raggedright ความเร่งโน้มถ่วง
\end{minipage}
\end{enumerate}
\vspace{-0.3\baselineskip}
\begin{small}
\textit{คำตอบ: ข}
\end{small}

\item ตัวต้านทาน 2 ตัวต่ออนุกรม 4 Ω และ 6 Ω ความต้านทานรวมเท่ากับ?
\begin{enumerate}[label=\alph*)., itemsep=0.3\baselineskip, topsep=0.3\baselineskip]
\item ก\ 
\begin{minipage}[t]{0.9\linewidth}
\raggedright 10 Ω
\end{minipage}
\item ข\ 
\begin{minipage}[t]{0.9\linewidth}
\raggedright 2.4 Ω
\end{minipage}
\item ค\ 
\begin{minipage}[t]{0.9\linewidth}
\raggedright 24 Ω
\end{minipage}
\item ง\ 
\begin{minipage}[t]{0.9\linewidth}
\raggedright 5 Ω
\end{minipage}
\end{enumerate}
\vspace{-0.3\baselineskip}
\begin{small}
\textit{คำตอบ: ก}
\end{small}

\item กฎของสเนลล์ (Snell's Law) สำหรับการหักเห: n₁ sinθ₁ = n₂ sinθ₂ โดย n คือ...?
\begin{enumerate}[label=\alph*)., itemsep=0.3\baselineskip, topsep=0.3\baselineskip]
\item ก\ 
\begin{minipage}[t]{0.9\linewidth}
\raggedright จำนวนเส้น
\end{minipage}
\item ข\ 
\begin{minipage}[t]{0.9\linewidth}
\raggedright ดรรชนีหักเห
\end{minipage}
\item ค\ 
\begin{minipage}[t]{0.9\linewidth}
\raggedright ความถี่
\end{minipage}
\item ง\ 
\begin{minipage}[t]{0.9\linewidth}
\raggedright ความเร็ว
\end{minipage}
\end{enumerate}
\vspace{-0.3\baselineskip}
\begin{small}
\textit{คำตอบ: ข}
\end{small}

\item สูตรเลนส์บาง: 1/f = 1/u + 1/v โดย u, v, f คืออะไร?
\begin{enumerate}[label=\alph*)., itemsep=0.3\baselineskip, topsep=0.3\baselineskip]
\item ก\ 
\begin{minipage}[t]{0.9\linewidth}
\raggedright u = ระยะภาพ, v = ระยะวัตถุ, f = ความยาวโฟกัส
\end{minipage}
\item ข\ 
\begin{minipage}[t]{0.9\linewidth}
\raggedright u = ระยะวัตถุ, v = ระยะภาพ, f = ความยาวโฟกัส
\end{minipage}
\item ค\ 
\begin{minipage}[t]{0.9\linewidth}
\raggedright u = ความสูงวัตถุ, v = ความสูงภาพ, f = กำลังขยาย
\end{minipage}
\item ง\ 
\begin{minipage}[t]{0.9\linewidth}
\raggedright u = ความยาวโฟกัส, v = ระยะวัตถุ, f = ระยะภาพ
\end{minipage}
\end{enumerate}
\vspace{-0.3\baselineskip}
\begin{small}
\textit{คำตอบ: ข}
\end{small}

\item สเปกตรัมเส้นของไฮโดรเจนที่เกิดจากอิเล็กตรอนตกจากระดับสูงลง n=2 เรียกว่าอะไร?
\begin{enumerate}[label=\alph*)., itemsep=0.3\baselineskip, topsep=0.3\baselineskip]
\item ก\ 
\begin{minipage}[t]{0.9\linewidth}
\raggedright Lyman series
\end{minipage}
\item ข\ 
\begin{minipage}[t]{0.9\linewidth}
\raggedright Balmer series
\end{minipage}
\item ค\ 
\begin{minipage}[t]{0.9\linewidth}
\raggedright Paschen series
\end{minipage}
\item ง\ 
\begin{minipage}[t]{0.9\linewidth}
\raggedright Brackett series
\end{minipage}
\end{enumerate}
\vspace{-0.3\baselineskip}
\begin{small}
\textit{คำตอบ: ข}
\end{small}

\item อนุภาคแอลฟา (α) มีองค์ประกอบและประจุอย่างไร?
\begin{enumerate}[label=\alph*)., itemsep=0.3\baselineskip, topsep=0.3\baselineskip]
\item ก\ 
\begin{minipage}[t]{0.9\linewidth}
\raggedright 2 โปรตอน + 2 นิวตรอน, ประจุ +2e
\end{minipage}
\item ข\ 
\begin{minipage}[t]{0.9\linewidth}
\raggedright 1 โปรตอน + 1 นิวตรอน, ประจุ +1e
\end{minipage}
\item ค\ 
\begin{minipage}[t]{0.9\linewidth}
\raggedright 2 อิเล็กตรอน, ประจุ -2e
\end{minipage}
\item ง\ 
\begin{minipage}[t]{0.9\linewidth}
\raggedright 1 อิเล็กตรอน, ประจุ -1e
\end{minipage}
\end{enumerate}
\vspace{-0.3\baselineskip}
\begin{small}
\textit{คำตอบ: ก}
\end{small}

\item ในการชนแบบยืดหยุ่น (Elastic collision) ข้อใดคงที่?
\begin{enumerate}[label=\alph*)., itemsep=0.3\baselineskip, topsep=0.3\baselineskip]
\item ก\ 
\begin{minipage}[t]{0.9\linewidth}
\raggedright เฉพาะพลังงานจลน์รวม
\end{minipage}
\item ข\ 
\begin{minipage}[t]{0.9\linewidth}
\raggedright เฉพาะโมเมนตัมรวม
\end{minipage}
\item ค\ 
\begin{minipage}[t]{0.9\linewidth}
\raggedright ทั้งพลังงานจลน์รวมและโมเมนตัมรวม
\end{minipage}
\item ง\ 
\begin{minipage}[t]{0.9\linewidth}
\raggedright ไม่มีอะไรคงที่
\end{minipage}
\end{enumerate}
\vspace{-0.3\baselineskip}
\begin{small}
\textit{คำตอบ: ค}
\end{small}

\item ศักย์ไฟฟ้า (Electric potential) V จากประจุ Q ที่ระยะ r มีค่าเท่าใด?
\begin{enumerate}[label=\alph*)., itemsep=0.3\baselineskip, topsep=0.3\baselineskip]
\item ก\ 
\begin{minipage}[t]{0.9\linewidth}
\raggedright V = kQ/r²
\end{minipage}
\item ข\ 
\begin{minipage}[t]{0.9\linewidth}
\raggedright V = kQ/r
\end{minipage}
\item ค\ 
\begin{minipage}[t]{0.9\linewidth}
\raggedright V = kQ²/r
\end{minipage}
\item ง\ 
\begin{minipage}[t]{0.9\linewidth}
\raggedright V = k/r
\end{minipage}
\end{enumerate}
\vspace{-0.3\baselineskip}
\begin{small}
\textit{คำตอบ: ข}
\end{small}

\item คลื่นแม่เหล็กไฟฟ้าใดมีความยาวคลื่นสั้นที่สุด?
\begin{enumerate}[label=\alph*)., itemsep=0.3\baselineskip, topsep=0.3\baselineskip]
\item ก\ 
\begin{minipage}[t]{0.9\linewidth}
\raggedright คลื่นวิทยุ
\end{minipage}
\item ข\ 
\begin{minipage}[t]{0.9\linewidth}
\raggedright รังสีอัลตราไวโอเลต
\end{minipage}
\item ค\ 
\begin{minipage}[t]{0.9\linewidth}
\raggedright รังสีแกมมา
\end{minipage}
\item ง\ 
\begin{minipage}[t]{0.9\linewidth}
\raggedright รังสีอินฟราเรด
\end{minipage}
\end{enumerate}
\vspace{-0.3\baselineskip}
\begin{small}
\textit{คำตอบ: ค}
\end{small}

\item การสลายตัวแบบ β⁻ เกิดจากการเปลี่ยนนิวตรอนเป็น...?
\begin{enumerate}[label=\alph*)., itemsep=0.3\baselineskip, topsep=0.3\baselineskip]
\item ก\ 
\begin{minipage}[t]{0.9\linewidth}
\raggedright โปรตอน + อิเล็กตรอน + antineutrino
\end{minipage}
\item ข\ 
\begin{minipage}[t]{0.9\linewidth}
\raggedright โปรตอนเท่านั้น
\end{minipage}
\item ค\ 
\begin{minipage}[t]{0.9\linewidth}
\raggedright อนุภาคแอลฟา
\end{minipage}
\item ง\ 
\begin{minipage}[t]{0.9\linewidth}
\raggedright รังสีแกมมาเท่านั้น
\end{minipage}
\end{enumerate}
\vspace{-0.3\baselineskip}
\begin{small}
\textit{คำตอบ: ก}
\end{small}

\item มวลวิกฤต (Critical mass) ของยูเรเนียม-235 ทรงกลมบริสุทธิ์มีค่าประมาณเท่าใด?
\begin{enumerate}[label=\alph*)., itemsep=0.3\baselineskip, topsep=0.3\baselineskip]
\item ก\ 
\begin{minipage}[t]{0.9\linewidth}
\raggedright 7 kg
\end{minipage}
\item ข\ 
\begin{minipage}[t]{0.9\linewidth}
\raggedright 52 kg
\end{minipage}
\item ค\ 
\begin{minipage}[t]{0.9\linewidth}
\raggedright 100 kg
\end{minipage}
\item ง\ 
\begin{minipage}[t]{0.9\linewidth}
\raggedright 235 kg
\end{minipage}
\end{enumerate}
\vspace{-0.3\baselineskip}
\begin{small}
\textit{คำตอบ: ข}
\end{small}

\item โมเมนตัมเชิงมุม L = Iω โดย I คือ... ?
\begin{enumerate}[label=\alph*)., itemsep=0.3\baselineskip, topsep=0.3\baselineskip]
\item ก\ 
\begin{minipage}[t]{0.9\linewidth}
\raggedright ความเร็วเชิงมุม
\end{minipage}
\item ข\ 
\begin{minipage}[t]{0.9\linewidth}
\raggedright แรงบิด
\end{minipage}
\item ค\ 
\begin{minipage}[t]{0.9\linewidth}
\raggedright โมเมนต์ความเฉื่อย
\end{minipage}
\item ง\ 
\begin{minipage}[t]{0.9\linewidth}
\raggedright คาบการหมุน
\end{minipage}
\end{enumerate}
\vspace{-0.3\baselineskip}
\begin{small}
\textit{คำตอบ: ค}
\end{small}

\item กฎของฟาราเดย์ (Faraday's Law) สำหรับการเหนี่ยวนำไฟฟ้า: EMF = -dΦ/dt หมายความว่าอะไร?
\begin{enumerate}[label=\alph*)., itemsep=0.3\baselineskip, topsep=0.3\baselineskip]
\item ก\ 
\begin{minipage}[t]{0.9\linewidth}
\raggedright แรงเคลื่อนไฟฟ้าเท่ากับอัตราการเปลี่ยนแปลงฟลักซ์แม่เหล็ก
\end{minipage}
\item ข\ 
\begin{minipage}[t]{0.9\linewidth}
\raggedright แรงเคลื่อนไฟฟ้าเท่ากับฟลักซ์แม่เหล็ก
\end{minipage}
\item ค\ 
\begin{minipage}[t]{0.9\linewidth}
\raggedright กระแสเท่ากับฟลักซ์
\end{minipage}
\item ง\ 
\begin{minipage}[t]{0.9\linewidth}
\raggedright สนามแม่เหล็กเท่ากับฟลักซ์
\end{minipage}
\end{enumerate}
\vspace{-0.3\baselineskip}
\begin{small}
\textit{คำตอบ: ก}
\end{small}

\item สมการเกรตติง: d sinθ = nλ โดย d คือ...?
\begin{enumerate}[label=\alph*)., itemsep=0.3\baselineskip, topsep=0.3\baselineskip]
\item ก\ 
\begin{minipage}[t]{0.9\linewidth}
\raggedright ระยะห่างระหว่างช่องกับจอ
\end{minipage}
\item ข\ 
\begin{minipage}[t]{0.9\linewidth}
\raggedright ระยะห่างระหว่างช่อง
\end{minipage}
\item ค\ 
\begin{minipage}[t]{0.9\linewidth}
\raggedright ความยาวคลื่น
\end{minipage}
\item ง\ 
\begin{minipage}[t]{0.9\linewidth}
\raggedright จำนวนช่อง
\end{minipage}
\end{enumerate}
\vspace{-0.3\baselineskip}
\begin{small}
\textit{คำตอบ: ข}
\end{small}

\item สมการของไอน์สไตน์สำหรับปรากฏการณ์โฟโตอิเล็กตริก: KE\_max = hf - φ ค่า φ คือ...?
\begin{enumerate}[label=\alph*)., itemsep=0.3\baselineskip, topsep=0.3\baselineskip]
\item ก\ 
\begin{minipage}[t]{0.9\linewidth}
\raggedright ความถี่ของแสง
\end{minipage}
\item ข\ 
\begin{minipage}[t]{0.9\linewidth}
\raggedright ฟังก์ชันงาน (Work function) ของโลหะ
\end{minipage}
\item ค\ 
\begin{minipage}[t]{0.9\linewidth}
\raggedright พลังงานจลน์สูงสุด
\end{minipage}
\item ง\ 
\begin{minipage}[t]{0.9\linewidth}
\raggedright ความยาวคลื่นของแสง
\end{minipage}
\end{enumerate}
\vspace{-0.3\baselineskip}
\begin{small}
\textit{คำตอบ: ข}
\end{small}

\item ถ้าพลังงานจากฟิชชัน 200 MeV ต่อปฏิกิริยา และโรงไฟฟ้า 1,000 MW ใช้เชื้อเพลิง 25,000 kg ต่อปี ประสิทธิภาพการใช้เชื้อเพลิงเป็นเท่าใด?
\begin{enumerate}[label=\alph*)., itemsep=0.3\baselineskip, topsep=0.3\baselineskip]
\item ก\ 
\begin{minipage}[t]{0.9\linewidth}
\raggedright 0.05\%
\end{minipage}
\item ข\ 
\begin{minipage}[t]{0.9\linewidth}
\raggedright 0.5\%
\end{minipage}
\item ค\ 
\begin{minipage}[t]{0.9\linewidth}
\raggedright 5\%
\end{minipage}
\item ง\ 
\begin{minipage}[t]{0.9\linewidth}
\raggedright 50\%
\end{minipage}
\end{enumerate}
\vspace{-0.3\baselineskip}
\begin{small}
\textit{คำตอบ: ข}
\end{small}

\item กฎแรงโน้มถ่วงของนิวตัน: F = Gm₁m₂/r² โดย G มีค่าเท่าใด?
\begin{enumerate}[label=\alph*)., itemsep=0.3\baselineskip, topsep=0.3\baselineskip]
\item ก\ 
\begin{minipage}[t]{0.9\linewidth}
\raggedright 6.67 × 10⁻¹¹ N·m²/kg²
\end{minipage}
\item ข\ 
\begin{minipage}[t]{0.9\linewidth}
\raggedright 9.8 N·m²/kg²
\end{minipage}
\item ค\ 
\begin{minipage}[t]{0.9\linewidth}
\raggedright 6.67 × 10¹¹ N·m²/kg²
\end{minipage}
\item ง\ 
\begin{minipage}[t]{0.9\linewidth}
\raggedright 6.67 × 10⁻¹¹ N·m/kg²
\end{minipage}
\end{enumerate}
\vspace{-0.3\baselineskip}
\begin{small}
\textit{คำตอบ: ก}
\end{small}

\item ในวงจร RC ค่าคงที่เวลา (Time constant) τ = RC หมายความว่าอะไร?
\begin{enumerate}[label=\alph*)., itemsep=0.3\baselineskip, topsep=0.3\baselineskip]
\item ก\ 
\begin{minipage}[t]{0.9\linewidth}
\raggedright เวลาที่ประจุเพิ่มขึ้นเป็น e เท่าของค่าเริ่มต้น
\end{minipage}
\item ข\ 
\begin{minipage}[t]{0.9\linewidth}
\raggedright เวลาที่ประจุลดลงเหลือ 1/e ของค่าเดิม
\end{minipage}
\item ค\ 
\begin{minipage}[t]{0.9\linewidth}
\raggedright ทั้ง ก และ ข
\end{minipage}
\item ง\ 
\begin{minipage}[t]{0.9\linewidth}
\raggedright เวลาที่กระแสเป็นศูนย์
\end{minipage}
\end{enumerate}
\vspace{-0.3\baselineskip}
\begin{small}
\textit{คำตอบ: ค}
\end{small}

\item คลื่นนิ่งในท่อที่ปลายปิดทั้งสองข้าง ความยาวท่อ L = n(λ/2) ค่า n คือ...?
\begin{enumerate}[label=\alph*)., itemsep=0.3\baselineskip, topsep=0.3\baselineskip]
\item ก\ 
\begin{minipage}[t]{0.9\linewidth}
\raggedright จำนวนเต็มบวกใดๆ
\end{minipage}
\item ข\ 
\begin{minipage}[t]{0.9\linewidth}
\raggedright จำนวนเต็มคี่
\end{minipage}
\item ค\ 
\begin{minipage}[t]{0.9\linewidth}
\raggedright จำนวนเต็มบวกที่ไม่เป็นศูนย์
\end{minipage}
\item ง\ 
\begin{minipage}[t]{0.9\linewidth}
\raggedright จำนวนเต็มคี่เท่านั้น
\end{minipage}
\end{enumerate}
\vspace{-0.3\baselineskip}
\begin{small}
\textit{คำตอบ: ก}
\end{small}

\item สารกัมมันตรังสีมีจำนวนนิวเคลียสเริ่มต้น 10¹² นิวเคลียส ค่าครึ่งชีวิต 5 ปี จะเหลือกี่นิวเคลียสหลัง 20 ปี?
\begin{enumerate}[label=\alph*)., itemsep=0.3\baselineskip, topsep=0.3\baselineskip]
\item ก\ 
\begin{minipage}[t]{0.9\linewidth}
\raggedright 6.25 × 10¹⁰
\end{minipage}
\item ข\ 
\begin{minipage}[t]{0.9\linewidth}
\raggedright 1.25 × 10¹¹
\end{minipage}
\item ค\ 
\begin{minipage}[t]{0.9\linewidth}
\raggedright 2.5 × 10¹⁰
\end{minipage}
\item ง\ 
\begin{minipage}[t]{0.9\linewidth}
\raggedright 5 × 10¹⁰
\end{minipage}
\end{enumerate}
\vspace{-0.3\baselineskip}
\begin{small}
\textit{คำตอบ: ก}
\end{small}

\item วัตถุถูกขว้างด้วยความเร็วต้น 30 m/s ที่มุม 30° จากพื้นราบ ความเร็วในแนวราบเริ่มต้นเท่าใด?
\begin{enumerate}[label=\alph*)., itemsep=0.3\baselineskip, topsep=0.3\baselineskip]
\item ก\ 
\begin{minipage}[t]{0.9\linewidth}
\raggedright 15 m/s
\end{minipage}
\item ข\ 
\begin{minipage}[t]{0.9\linewidth}
\raggedright 26 m/s
\end{minipage}
\item ค\ 
\begin{minipage}[t]{0.9\linewidth}
\raggedright 30 m/s
\end{minipage}
\item ง\ 
\begin{minipage}[t]{0.9\linewidth}
\raggedright 7.5 m/s
\end{minipage}
\end{enumerate}
\vspace{-0.3\baselineskip}
\begin{small}
\textit{คำตอบ: ข}
\end{small}

\item หลอดไฟ 100 W ต่อกับไฟบ้าน 220 V จะมีกระแสไฟฟ้าผ่านเท่าใด?
\begin{enumerate}[label=\alph*)., itemsep=0.3\baselineskip, topsep=0.3\baselineskip]
\item ก\ 
\begin{minipage}[t]{0.9\linewidth}
\raggedright 0.45 A
\end{minipage}
\item ข\ 
\begin{minipage}[t]{0.9\linewidth}
\raggedright 2.2 A
\end{minipage}
\item ค\ 
\begin{minipage}[t]{0.9\linewidth}
\raggedright 4.5 A
\end{minipage}
\item ง\ 
\begin{minipage}[t]{0.9\linewidth}
\raggedright 0.22 A
\end{minipage}
\end{enumerate}
\vspace{-0.3\baselineskip}
\begin{small}
\textit{คำตอบ: ก}
\end{small}

\item รถยนต์เคลื่อนที่เข้าหาผู้สังเกตที่หยุดนิ่งด้วยความเร็ว 30 m/s ความถี่เสียงจริง 500 Hz ความเร็วเสียง 340 m/s ความถี่ที่ได้ยินเป็นเท่าใด?
\begin{enumerate}[label=\alph*)., itemsep=0.3\baselineskip, topsep=0.3\baselineskip]
\item ก\ 
\begin{minipage}[t]{0.9\linewidth}
\raggedright 500 Hz
\end{minipage}
\item ข\ 
\begin{minipage}[t]{0.9\linewidth}
\raggedright 544 Hz
\end{minipage}
\item ค\ 
\begin{minipage}[t]{0.9\linewidth}
\raggedright 456 Hz
\end{minipage}
\item ง\ 
\begin{minipage}[t]{0.9\linewidth}
\raggedright 585 Hz
\end{minipage}
\end{enumerate}
\vspace{-0.3\baselineskip}
\begin{small}
\textit{คำตอบ: ข}
\end{small}

\item ในข้อสอบนิวเคลียร์ เมื่อโจทย์ให้มวลนิวเคลียสและถามพลังงานยึดเหนี่ยว ต้องระวังอะไรเป็นพิเศษ?
\begin{enumerate}[label=\alph*)., itemsep=0.3\baselineskip, topsep=0.3\baselineskip]
\item ก\ 
\begin{minipage}[t]{0.9\linewidth}
\raggedright ต้องใช้หน่วย kg เท่านั้น
\end{minipage}
\item ข\ 
\begin{minipage}[t]{0.9\linewidth}
\raggedright ต้องระวังโจทย์ให้มวลอะตอมหรือมวลนิวเคลียส เพราะต่างกันที่มวลอิเล็กตรอน
\end{minipage}
\item ค\ 
\begin{minipage}[t]{0.9\linewidth}
\raggedright ต้องใช้สมการ E = mc
\end{minipage}
\item ง\ 
\begin{minipage}[t]{0.9\linewidth}
\raggedright ไม่ต้องระวังอะไร
\end{minipage}
\end{enumerate}
\vspace{-0.3\baselineskip}
\begin{small}
\textit{คำตอบ: ข}
\end{small}

\item เทคนิค 'ตัดตัวเลือก' ที่มีประสิทธิภาพสูงในข้อสอบ PAT3 คืออะไร?
\begin{enumerate}[label=\alph*)., itemsep=0.3\baselineskip, topsep=0.3\baselineskip]
\item ก\ 
\begin{minipage}[t]{0.9\linewidth}
\raggedright ตัดตัวเลือกที่มีตัวเลขสูงสุดเสมอ
\end{minipage}
\item ข\ 
\begin{minipage}[t]{0.9\linewidth}
\raggedright ตัดตัวเลือกที่มีหน่วยผิดออกก่อน แล้วตรวจสอบทิศทางของปริมาณ (บวก/ลบ, ใหญ่/เล็ก)
\end{minipage}
\item ค\ 
\begin{minipage}[t]{0.9\linewidth}
\raggedright ตัดตัวเลือก ก และ ข ออกเสมอ
\end{minipage}
\item ง\ 
\begin{minipage}[t]{0.9\linewidth}
\raggedright ตัดตัวเลือกที่มี π ออก
\end{minipage}
\end{enumerate}
\vspace{-0.3\baselineskip}
\begin{small}
\textit{คำตอบ: ข}
\end{small}

\end{enumerate}
\end{document}